\section{Clases y Estrategias de Estudio}

\subsection{Clases Introductorias}
Inician generalmente una o dos semanas antes del ciclo regular (aproximadamente el 2 de marzo). No son obligatorias ni afectan tu promedio, pero son vitales para no "volar" la primera semana.

\subsection{Libros y Consejos}
La universidad exige ser autodidacta. Aquí algunos tips de supervivencia:
\begin{itemize}
    \item \textbf{Bibliografía base:} Para Física usa el \textit{Serway}. Para Geometría se recomienda fuertemente el libro \textit{Vidal} (excelente para ponderar) y el \textit{Nakamura}. 
    \item \textbf{Material FIM:} En las fotocopiadoras de la facultad venden cuadernos de cursos específicos escritos por los mismos ingenieros. Consigue el cuaderno de ese amigo que siempre asiste a clases; los profesores suelen resolver ejercicios muy parecidos a los del examen días antes de la prueba.
    \item \textbf{Laboratorios:} Busca armar un buen grupo de laboratorio desde el día uno, te salvará la vida. Apóyate mucho en YouTube.
\end{itemize}
